\section{Method}
\label{sec:method}

\Cref{fig:architecture} lays out the two-path world model
(\method) end to end. This section defines it precisely: the shared
encoder and two prediction heads (\Cref{sec:method-forward}), the
objective that couples them (\Cref{sec:method-loss}), and the four
structural properties that make the recovered physical parameters a
claim about the model rather than an artifact of how it was trained
(\Cref{sec:method-invariants}).

% Two-path architecture figure. Pure TikZ -- every box/label is real text,
% not a raster image.
%
% Topology matches stable_worldmodel/wm/physwm/physwm.py exactly: encoder
% and predictor are a SINGLE shared trunk producing one action-conditioned
% latent z_hat, which forks into the two heads. Loss formulas are NOT drawn
% here -- they are Eq. (4) in the text; the figure's only job is what feeds
% what. Sized to its natural extent, not stretched to fill the column.
\begin{figure*}[t]
\centering
\resizebox{\textwidth}{!}{%
\begin{tikzpicture}[
  node distance=7mm and 9mm,
  box/.style={draw, rounded corners=2pt, align=center, font=\scriptsize,
              minimum height=7mm, inner sep=3pt},
  learned/.style={box, fill=blue!7, draw=blue!55},
  physical/.style={box, fill=orange!8, draw=orange!65!black},
  shared/.style={box, fill=gray!12, draw=gray!60, thick},
  frozen/.style={box, fill=gray!25, draw=black, thick, dashed},
  gt/.style={box, fill=green!8, draw=green!45!black},
  lossdot/.style={draw, circle, fill=red!8, draw=red!55!black,
              font=\scriptsize\itshape, inner sep=1pt, minimum size=7mm},
  arr/.style={-{Stealth[length=1.8mm]}, thick},
  detacharr/.style={-{Stealth[length=1.8mm]}, thick, dash pattern=on 2pt off 2pt, red!55!black},
  ablarr/.style={-{Stealth[length=1.5mm]}, thin, dotted, gray!70},
]

% --- shared trunk: encoder and predictor produce ONE action-conditioned
% latent z_hat, read by both heads ---
\node[box, fill=white] (pixels) {pixels / state\\ $o_{\le t}$};
\node[shared, right=of pixels] (encoder) {encoder\\ $z=f_\phi(o_{\le t})$};
\node[shared, right=of encoder, minimum width=6mm] (z) {$z$};
\node[shared, right=of z, text width=17mm] (predictor) {predictor\\ block-causal};
\node[shared, right=of predictor, minimum width=7mm] (zhat) {$\hat z$};

% --- Path A: learned, upper branch off z_hat ---
\node[learned, above right=8mm and 9mm of zhat] (decoder) {state decoder};
\node[learned, right=of decoder, minimum width=6mm] (stateA) {$s_A$};
\node[gt, above=5mm of stateA] (target) {dataset $s_{\mathrm{next}}$};
\node[lossdot, right=of stateA] (LA) {$\mathcal L_A$};

% --- Path B: physical, lower branch off z_hat ---
\node[physical, below right=8mm and 9mm of zhat] (probe) {probe (linear)};
\node[physical, right=of probe, minimum width=6mm] (theta) {$\theta$};
\node[frozen, right=of theta, text width=18mm] (solver) {frozen solver\\ $\dot v = F(s,a,\theta)/m$};
\node[physical, right=of solver, minimum width=6mm] (stateB) {$s_B$};
\node[box, fill=white, below=5mm of probe] (sa) {$(s_t,a_t)$};
\node[lossdot, right=of stateB] (LB) {$\mathcal L_B$};

% --- trunk wiring: strictly sequential, no fork until z_hat ---
\draw[arr] (pixels) -- (encoder);
\draw[arr] (encoder) -- (z);
\draw[arr] (z) -- (predictor);
\draw[arr] (predictor) -- (zhat);
\draw[arr] (zhat.north) |- (decoder.west);
\draw[arr] (zhat.south) |- (probe.west);

% --- Path A wiring ---
\draw[arr] (decoder) -- (stateA);
\draw[arr] (stateA) -- (LA);
\draw[arr] (target) -- (LA);

% --- Path B wiring ---
\draw[arr] (probe) -- (theta);
\draw[arr] (theta) -- (solver);
\draw[arr] (sa) -- (solver);
\draw[arr] (solver) -- (stateB);
\draw[arr] (stateB) -- (LB);

% --- the key edge: Path A's own prediction as Path B's target ---
\draw[detacharr] (stateA.south) .. controls +(0,-9mm) and +(0,9mm) .. (LB.north);
\node[red!55!black, font=\tiny\itshape, align=left]
  at ($(stateA)!0.5!(LB)+(9mm,0)$) {stop-gradient\\ target};

% --- gradient path back into probe, through the frozen solver ---
\draw[arr, orange!65!black, -{Stealth[length=1.6mm]}, dashed]
  (LB.south) -- ++(0,-3mm) -| (solver.south);
\node[orange!65!black, font=\tiny\itshape] at ($(solver)+(0,-11mm)$)
  {$\nabla$ through the frozen solver};

% --- the ablation this topology rules out ---
\draw[ablarr] (z.south) .. controls +(0,-13mm) and +(-8mm,0) .. (probe.west);
\node[gray!60, font=\tiny\itshape] at ($(z)+(-1mm,-12mm)$) {ablation: probe reads $z$};

\end{tikzpicture}%
}
\caption{\textbf{\method architecture.} Encoder and predictor form one
shared trunk producing a single action-conditioned latent $\hat z$; both
heads read $\hat z$, never a private view of the observation. \pathA
(blue) is the ordinary predictor: a decoder reads $\hat z$ to a state
$s_A$, and $\mathcal L_A$ (\Cref{eq:loss}) regresses the dataset's
ground truth, exactly what a predictive world model already does.
\pathB (orange) reads a physics vector $\theta$ off the same $\hat z$
with a deliberately low-capacity probe, integrates it through a
\emph{frozen, zero-parameter} solver (gray, dashed) to get an
independent estimate $s_B$. $\mathcal L_B$ does \emph{not} use the
dataset label: it regresses $s_A$ with gradient stopped (red, dashed),
so the only route back to a learnable parameter is through the fixed
solver into $\theta$ (orange dashed arrow), and $\theta$ can only
improve by becoming a better physical explanation of what \pathA
already predicts. The dotted arrow is not part of the method: it marks
the pre-action routing ablation (\Cref{sec:setup-benchmarks}), where the
probe reads $z$ instead of $\hat z$.}
\label{fig:architecture}
\end{figure*}


\subsection{Shared latent, two heads}
\label{sec:method-forward}

Given an observation window $o_{1:T}$ (pixels, or raw state for the
ablation in \Cref{sec:experimental-setup}), a single encoder $f_\phi$
produces per-frame patch tokens $z_t = f_\phi(o_t) \in \R^{N \times
D}$. Neither head has a private view of the observation: both are
computed from $z$, and, as the rest of this section makes precise, both
read it through the same predictor, never one reading $z$ directly
while the other reads the predictor's output.

\pathA is a standard learned predictor. A block-causal
spatio-temporal transformer $g_\psi$ \citep{zhou2024dinowm} attends
over all patches of frames $\le t$ and predicts a residual update to
the next frame's tokens,
\begin{equation}
  \hat z_{t+1} = z_t + g_\psi(z_{1:t}, a_{1:t}),
\end{equation}
and a decoder $d_\chi$ pools $\hat z_{t+1}$ over patches and reads out
a state estimate $s_{A,t+1} = d_\chi(\hat z_{t+1})$ (\stateA{} once the
timestep is unambiguous, as in \Cref{eq:loss}). Nothing here
is physics-specific; \pathA's only job is to be an accurate,
general-purpose forecaster, trained the way any DINO-WM-style
predictor is trained.

\pathB is a physical predictor built from a second, deliberately
low-capacity probe. $h_\rho$ reads the \emph{same} action-conditioned
tokens $\hat z_{1:t}$ that \pathA's decoder reads, not the raw encoder
output $z$, restricted to the frames the solver is about to step from,
never the frame being predicted, so the probe cannot shortcut by seeing
the answer. Reading $z$ directly, before the predictor conditions it on
the action, is kept only as a falsifying ablation
(\Cref{sec:setup-benchmarks}): if it scores the same as reading
$\hat z$, then sharing the predictor's action-conditioned latent between
the two paths is decoration, not a load-bearing design choice. The probe
emits an unconstrained parameter vector, which we squash into a fixed,
physically sane range with a sigmoid reparameterization:
\begin{equation}
  \theta_{\mathrm{raw}} = h_\rho(\hat z_{1:t}), \qquad
  \theta = \mathrm{lo} + (\mathrm{hi} - \mathrm{lo}) \odot
    \sigma(\theta_{\mathrm{raw}}) \in [\mathrm{lo}, \mathrm{hi}]^K.
\end{equation}
$\mathrm{lo}, \mathrm{hi} \in \R^K$ are fixed constants, not learned.
By default $h_\rho$ is a single linear map and $\theta$ is one vector
per episode, pooled over patches and over the whole context window,
rather than one per timestep. Physical parameters are properties of
the system, not of the instant, and a probe with almost no capacity is
the honest test of whether $\theta$ is linearly present in $z$ at all.
$\theta$ is then handed, with the current physical state and action,
to a fixed integrator $S$,
\begin{equation}
  s_{B,t+1} = S(s_t, a_t, \theta),
\end{equation}
which produces an independent, physically-derived next-state estimate.
$S$ is hand-derived per benchmark (a point-mass contact-and-drag
model, a rigid cart-pole, a quasi-static planar-pushing model;
\Cref{sec:experimental-setup} gives the exact equations), and it is
never trained. $S$ owns zero learnable parameters, so every degree of
freedom \pathB has for explaining a transition lives in $\theta$.

\subsection{Objective: the model as its own supervisor}
\label{sec:method-loss}

All losses are computed in a fixed, dataset-fitted normalized state
space (z-scored per channel, held in buffers rather than learned, so a
checkpoint is self-contained), which keeps $\mathcal L_A$ and
$\mathcal L_B$ commensurate and $\alpha,\beta$ scale-free:
\begin{equation}
  \label{eq:loss}
  \mathcal L \;=\; \underbrace{\big\lVert \stateA - \stategt
    \big\rVert^2}_{\mathcal L_A}
  \;+\; \alpha \underbrace{\big\lVert \stateB - \stopgrad{\stateA}
    \big\rVert^2}_{\mathcal L_B}
  \;+\; \beta \underbrace{\big\lVert \stateA - \stateB \big\rVert^2}
    _{\mathcal L_{\mathrm{consistency}}},
\end{equation}
with $\beta = 0$ by default. Everything this paper is built around is
in the target of $\mathcal L_B$: $\stopgrad{\stateA}$, \pathA's own
prediction with gradient stopped, in place of the dataset label
$\stategt$.

Once both paths exist, three candidate targets for $\mathcal L_B$ are
available. Comparing $\stateB$ to itself is vacuous; nothing
constrains $\theta$. Comparing it to $\stategt$ is the obvious first
choice. Comparing it to $\stateA$ is what we use. We reject the second
option for a specific reason: fitting the physics path directly to the
raw benchmark label tests whether $S(\cdot,\cdot,\theta)$ can be made
to match a transition for some $\theta$, a question $\mathcal L_A$
already answers with a higher-capacity function class and, typically,
lower error. It does not test whether $\theta$ explains anything the
model itself relies on. Using $\stopgrad{\stateA}$ instead asks the
sharper question: starting from a prediction the model already trusts,
can a near-zero-capacity, physically-constrained probe recover the
compact parameterization underneath it? That is the top-down step from
\Cref{sec:introduction}. The reliable upper level of the
representation, $\stateA$, becomes the training signal for the thin
lower level, $\theta$, and \Cref{eq:loss} never uses a
physical-parameter annotation to do it.

The stop-gradient is what keeps this well-posed rather than circular.
$\mathcal L_B$ never receives gradient through $\stateA$, so it cannot
move \pathA's decoder or predictor toward the solver's output; the
pressure runs one way, from the general prediction onto the physics
probe. Drop the stop-gradient and minimizing \Cref{eq:loss} jointly
admits a degenerate shortcut: \pathA drifts toward whatever \pathB
finds easiest to reproduce, which would make a bad $\theta$ look
consistent by pulling the general predictor down to meet it rather
than by explaining it.

\subsection{Design invariants}
\label{sec:method-invariants}

Four structural properties make the recovered $\theta$ worth testing.
Each is enforced as a checked invariant in the implementation, not
left as a training convention that could silently drift.

\begin{enumerate}[leftmargin=1.6em, itemsep=2pt]
  \item \textbf{One latent, two predictions, neither post-hoc.}
    $\stateA$ and $\stateB$ both branch from the same $z$
    (\Cref{sec:method-forward}) and are computed independently; they
    are compared only at the loss, so neither is fit to the other
    after the fact.
  \item \textbf{\pathA regresses ground truth; \pathB regresses
    \pathA; never the reverse.} $\mathcal L_A$ is the only term whose
    gradient into $d_\chi$ and $g_\psi$ originates from $\stategt$.
    $\mathcal L_B$'s target is detached (\Cref{sec:method-loss}), so
    with $\beta=0$ no configuration of \Cref{eq:loss} pushes
    $\mathcal L_B$'s gradient into \pathA.
  \item \textbf{$\theta$ is always a forward-pass output, never a free
    variable.} $\theta = \mathrm{bound\_theta}(h_\rho(z))$ on every
    forward pass, never an optimizable parameter fit directly to a
    target. If $\theta$ could be a free variable, $\mathcal L_B$ could
    be driven to zero by memorizing a per-episode value independent of
    $z$, and the experiment would say nothing about the latent.
  \item \textbf{$S$ is frozen and differentiable.} $S$ has zero
    learnable parameters, constants live in fixed buffers rather than
    \texttt{nn.Parameter}s, so gradient flows through its fixed
    equations of motion into $\theta$ and hence into $h_\rho$, but
    nothing inside $S$ ever changes. The only way to lower
    $\mathcal L_B$ is to move $\theta$ in a direction the fixed physics
    of $S$ turns into a better match for \pathA. Whether that direction
    is also a physically meaningful one is exactly what the
    identifiability and transfer experiments in
    \Cref{sec:experimental-setup} test.
\end{enumerate}

The optional consistency term $\mathcal L_{\mathrm{consistency}}$ in
\Cref{eq:loss} is symmetric between $\stateA$ and $\stateB$ and off by
default. It is an ablation knob, not part of the main recipe: a
nonzero, non-symmetric variant would reintroduce the coupling
invariant~2 rules out.
